\chapter{Categories}

\begin{overviewbox}
Chapter~1 introduced four ideas: objects, morphisms, composition, and identity.
Each came with rules --- composition is associative, identity is neutral --- but
we never packaged all four ideas and their rules into a single named thing.

This chapter does that. A \emph{category} is the name for the package. Once we
have the name and the definition, we can look at many different situations and
ask: ``is this a category?'' The answer is often yes --- and recognizing that
two very different-looking situations are both categories is the first payoff of
the theory.

The chapter also introduces \emph{isomorphism} --- the precise meaning of
``these two objects are the same for our purposes'' --- and the \emph{opposite
category}, which reverses all arrows and will recur throughout the book.

The Formal Restatement at the end collects the definitions in compact notation.
\end{overviewbox}


% ═══════════════════════════════════════════════════════════════════════════
\section{The Definition of a Category}
\label{sec:category-definition}
% ═══════════════════════════════════════════════════════════════════════════

\subsection{Informally: Packaging What We Have}

We have been working with objects and morphisms for an entire chapter. We have
two rules: composition is associative, and identity morphisms are neutral. Now
we bundle everything together and give the bundle a name.

A \term{category} is a collection of objects and morphisms, equipped with
composition and identity morphisms, satisfying the associativity and unit laws.

That's it. A category is not a complicated new idea. It is a name for the
structure we have already been using.

\subsection{Expanding: Checking Our Examples}

The three running examples from Chapter~1 are all categories. Let's verify
this for one of them in detail.

\begin{example}{The Recipe as a Category}
Recall the recipe with three steps: Chop, Sauté, Add.

\medskip
\begin{center}
\begin{tikzcd}[column sep=2.5em]
  \text{Chop}
    \arrow[r, "f"]
    \arrow[loop left, "\mathrm{id}_C"]
  & \text{Sauté}
    \arrow[r, "g"]
    \arrow[loop above, "\mathrm{id}_S"]
  & \text{Add}
    \arrow[loop right, "\mathrm{id}_A"]
\end{tikzcd}
\end{center}
\medskip

We check each requirement:

\begin{itemize}
  \item \textbf{Objects?} Yes: Chop, Sauté, Add.
  \item \textbf{Morphisms?} Yes: $f$, $g$, $g \circ f$, and the three
        identity morphisms.
  \item \textbf{Composition?} Yes: $g \circ f$ exists and goes from Chop
        to Add, as required.
  \item \textbf{Identity?} Yes: each step has a ``do nothing at this step''
        morphism.
  \item \textbf{Associativity?} There are only two non-identity, non-composed
        morphisms ($f$ and $g$), so the only chain of three is degenerate ---
        associativity holds vacuously.
  \item \textbf{Unit laws?} $f \circ \mathrm{id}_C = f$,\;
        $\mathrm{id}_S \circ f = f$. Checked.
\end{itemize}

The recipe is a category.
\end{example}

\begin{dialog}
When I want to verify that something is a category, I run through a
checklist:
\begin{enumerate}
  \item Can I name the objects?
  \item Can I name the morphisms, with domains and codomains?
  \item Is composition defined wherever it should be, and associative?
  \item Does every object have an identity morphism that is neutral?
\end{enumerate}
If all four answers are yes, I have a category. This checklist applies
to every example in this book and every example I'll encounter elsewhere.
\end{dialog}

\subsection{Formalizing: The Four-Part Definition}

A \term{category} $\mathcal{C}$ consists of four pieces of data:

\begin{enumerate}
  \item A collection $\mathrm{Ob}(\mathcal{C})$, whose elements are called
        \term{objects}.
  \item For each pair of objects $A, B \in \mathrm{Ob}(\mathcal{C})$, a
        collection $\mathrm{Hom}_{\mathcal{C}}(A, B)$, whose elements are
        called \term{morphisms} from $A$ to $B$.
  \item For each triple of objects $A, B, C$, a \term{composition} function
        \[
          \circ\;:\;\mathrm{Hom}(B, C) \times \mathrm{Hom}(A, B)
          \;\longrightarrow\; \mathrm{Hom}(A, C).
        \]
  \item For each object $A$, a distinguished \term{identity morphism}
        $\mathrm{id}_A \in \mathrm{Hom}(A, A)$.
\end{enumerate}

This data must satisfy two \term{axioms}:

\begin{itemize}
  \item \textbf{Associativity}: For all $f \in \mathrm{Hom}(A,B)$,
        $g \in \mathrm{Hom}(B,C)$, $h \in \mathrm{Hom}(C,D)$:
        \[
          (h \circ g) \circ f = h \circ (g \circ f).
        \]
  \item \textbf{Unit laws}: For all $f \in \mathrm{Hom}(A,B)$:
        \[
          f \circ \mathrm{id}_A = f = \mathrm{id}_B \circ f.
        \]
\end{itemize}

\begin{readernote}
When you meet a definition with this structure --- ``consists of'' followed
by ``satisfying'' --- read the ``consists of'' part as a list of
\emph{what you must supply}, and the ``satisfying'' part as a list of
\emph{what your supply must pass}. The data is what you bring to the table;
the axioms are the tests your data must survive.
\end{readernote}

\subsection{Reorganizing: The Examples Pass the Tests}

\textbf{Recipe category.} We supply: objects $\{\mathrm{Chop}, \mathrm{Sauté},
\mathrm{Add}\}$; morphisms $f$, $g$, $g \circ f$, and three identities;
composition by the rule of chaining steps; identity by the ``stay at this
step'' morphism. The associativity test passes (checked above). The unit
test passes.

\textbf{Transit category.} The stations are objects. The train lines are
morphisms. The ``stay at station'' routes are identities. Any journey along
lines in sequence is a composite. Since real transit routes are unambiguous
regardless of how you mentally group the legs, associativity holds. The
transit map is a category.

\textbf{Temperature ordering.} The readings $\{0°\text{C}, 20°\text{C},
37°\text{C}, 100°\text{C}\}$ are objects. A morphism $A \to B$ exists
exactly when $A \leq B$ (``is no warmer than''). Identity: every temperature
is no warmer than itself. Composition: if $A \leq B$ and $B \leq C$ then
$A \leq C$. Associativity holds because $\leq$ is transitive, and there is
at most one morphism between any two objects so associativity is automatic.
This is a category.

\subsection{Simplifying: The Definition in Brief}

A category packages four things: objects, morphisms, composition,
identity. It satisfies two laws:
\[
  (h \circ g) \circ f = h \circ (g \circ f)
  \qquad\text{and}\qquad
  f \circ \mathrm{id}_A = f = \mathrm{id}_B \circ f.
\]
Every example we have seen passes both laws. These two equations are the
entire content of the definition.

\subsection{Formally: The Standard Formulation}

\formalnote{Notation}
We write $\mathcal{C}$, $\mathcal{D}$, $\mathcal{E}$ for categories and
use bold letters $\mathbf{C}$, $\mathbf{D}$ in some contexts. The notation
$f \in \mathcal{C}(A, B)$ is an alternative to $f \in \mathrm{Hom}_\mathcal{C}(A,B)$.
When the category is clear from context we write $\mathrm{Hom}(A,B)$.

\formalnote{Size}
In the definition above we said ``collection'' rather than ``set.'' The
distinction matters: if $\mathrm{Ob}(\mathcal{C})$ is a set, $\mathcal{C}$
is called \term{small}; if the hom-sets $\mathrm{Hom}(A,B)$ are all sets,
$\mathcal{C}$ is called \term{locally small}. Most categories we encounter are
at least locally small. We will not dwell on set-theoretic foundations here;
the interested reader may consult Appendix~A (to be written).

\formalnote{Uniqueness of identities revisited}
The proof from Chapter~1 that identity morphisms are unique works
uniformly in any category: if $e, e'$ both satisfy the unit laws at $A$,
then $e = e \circ e' = e'$. No additional hypotheses are needed.


% ═══════════════════════════════════════════════════════════════════════════
\section{A Catalogue of Small Categories}
\label{sec:small-categories}
% ═══════════════════════════════════════════════════════════════════════════

\subsection{Informally: Categories Can Be Tiny}

Not all categories are large. Some have very few objects and morphisms.
These tiny categories are not toys --- they are the building blocks used
throughout the rest of category theory to state definitions and theorems
in a compact, diagram-based way.

We build five small categories from scratch.

\subsection{Expanding: Five Small Categories}

\begin{example}{$\mathbf{0}$: The Empty Category}
Objects: none. Morphisms: none.

Does it satisfy the axioms? Every axiom has the form ``for all objects
$A$\ldots'' or ``for all morphisms $f$\ldots''. With nothing to quantify over,
every such statement is true \emph{vacuously} --- it holds because there is
nothing that could fail it.

The empty category is valid. It is the categorical equivalent of a blank
page: a place where nothing has been written, not a place where something
is wrong.
\end{example}

\begin{example}{$\mathbf{1}$: One Object, One Morphism}
Objects: one object, call it $*$. Morphisms: one morphism, $\mathrm{id}_*$.

The only composition is $\mathrm{id}_* \circ \mathrm{id}_* = \mathrm{id}_*$.
The only unit law is $\mathrm{id}_* \circ \mathrm{id}_* = \mathrm{id}_*$.
Both are satisfied. The category $\mathbf{1}$ is the simplest non-empty
category: one dot, one loop.
\end{example}

\begin{example}{$\mathbf{2}$: Two Objects, One Arrow}
Objects: $0$ and $1$. Morphisms: $\mathrm{id}_0$, $\mathrm{id}_1$, and one
morphism $\iota : 0 \to 1$.

\medskip
\begin{center}
\begin{tikzcd}
  0 \arrow[loop left, "\mathrm{id}_0"] \arrow[r, "\iota"] &
  1 \arrow[loop right, "\mathrm{id}_1"]
\end{tikzcd}
\end{center}
\medskip

There is no morphism from $1$ back to $0$. The only non-trivial
composition would require a morphism starting at $1$, but there is none
(besides $\mathrm{id}_1$ which goes $1 \to 1$). So composition is
complete. The unit laws hold for $\iota$ because
$\iota \circ \mathrm{id}_0 = \iota = \mathrm{id}_1 \circ \iota$.

The category $\mathbf{2}$ is sometimes called the \term{walking morphism}:
it contains exactly one non-identity morphism and nothing else. It is used
later to ``pick out'' a single morphism from a category.
\end{example}

\begin{example}{An Ordering Category}
Take any collection of objects with a relation $\leq$ that is:
\begin{itemize}
  \item \textbf{Reflexive}: $A \leq A$ for all $A$ (gives identity),
  \item \textbf{Transitive}: $A \leq B$ and $B \leq C$ implies $A \leq C$
        (gives composition).
\end{itemize}

Define $\mathrm{Hom}(A, B)$ to contain exactly one morphism when $A \leq B$,
and to be empty otherwise. Since there is at most one morphism between any
two objects, associativity is automatic. Identity comes from reflexivity.

Such a category is called a \term{preorder category}. Our temperature
example is a preorder category on $\{0°\text{C}, 20°\text{C}, 37°\text{C},
100°\text{C}\}$ with the usual $\leq$ on temperatures.

A preorder category where $A \leq B$ and $B \leq A$ implies $A = B$ is
called a \term{poset category}.
\end{example}

\begin{example}{A One-Object Category}
Take a single object $*$. Let the morphisms from $*$ to $*$ be some
collection $M$ of ``actions'' that can be composed in sequence, where
composition is associative and there is a ``do nothing'' action.

For concreteness: suppose $*$ is a light panel, and the morphisms are
sequences of button presses. Pressing button A then button B is a composite
action. The ``do nothing'' action is the identity. If button presses compose
associatively (which they do, since doing A then B then C is the same
regardless of grouping), this is a category.

A one-object category is precisely a \term{monoid}: a set with an
associative binary operation and an identity element. We will return to
this connection when we discuss functors.
\end{example}

\begin{dialog}
When I encounter a small category, I draw it as a directed graph and check
two things:
\begin{enumerate}
  \item Is the graph \emph{closed under composition}? For every two
        composable edges, is there an edge representing their composite?
  \item Does every vertex have a \emph{loop} for its identity?
\end{enumerate}
If both hold and composition is associative, I have a category. For tiny
categories, associativity is often automatic because there simply are not
enough morphisms to chain in three different ways.
\end{dialog}

\subsection{Formalizing: What Makes a Category Small?}

In the example above we used the word ``small'' informally. Let's be
precise.

\begin{itemize}
  \item $\mathcal{C}$ is \term{small} if $\mathrm{Ob}(\mathcal{C})$ is a
        set (not a proper class).
  \item $\mathcal{C}$ is \term{finite} if $\mathrm{Ob}(\mathcal{C})$ and all
        hom-sets are finite sets.
  \item $\mathcal{C}$ is \term{discrete} if the only morphisms are
        identity morphisms: $\mathrm{Hom}(A,B) = \varnothing$ for $A \neq B$
        and $\mathrm{Hom}(A,A) = \{\mathrm{id}_A\}$.
  \item $\mathcal{C}$ is \term{thin} (or a \term{preorder}) if there is at
        most one morphism between any two objects.
\end{itemize}

The categories $\mathbf{0}$, $\mathbf{1}$, $\mathbf{2}$ are all small and
finite. The temperature example is small, finite, and thin. A discrete
category on $n$ objects has $n$ objects and $n$ morphisms (just identities).

\subsection{Reorganizing: Naming What We Already Know}

Our examples from Chapter~1 now have names:
\begin{itemize}
  \item The recipe (Chop $\to$ Sauté $\to$ Add) is a small, finite,
        \emph{thin} category --- there is exactly one morphism from any
        step to any step that comes later, and none in the other direction.
  \item The temperature ordering is a small, finite, \emph{poset} category.
  \item The approval chain is a small, finite, thin category, identical in
        shape to the recipe.
\end{itemize}

The single-object category $\mathbf{1}$ is both a category and a monoid.
We will see later that any monoid can be viewed as a one-object category,
and vice versa.

\subsection{Simplifying: A Table of Small Categories}

\medskip
\begin{center}
\begin{tabular}{lllll}
  Name & Objects & Non-id morphisms & Thin? & Finite? \\[4pt]
  \hline\\[-6pt]
  $\mathbf{0}$ & 0 & 0 & yes & yes \\
  $\mathbf{1}$ & 1 & 0 & yes & yes \\
  $\mathbf{2}$ & 2 & 1 & yes & yes \\
  Preorder & any & varies & yes & varies \\
  One-object (monoid) & 1 & varies & no & varies \\
\end{tabular}
\end{center}
\medskip

\subsection{Formally: Discrete and Indiscrete Categories}

\formalnote{Discrete category on a collection $S$}
Given any collection $S$, the \term{discrete category} $\mathrm{Disc}(S)$
has $\mathrm{Ob} = S$ and $\mathrm{Hom}(A,B) = \varnothing$ for $A \neq B$,
$\mathrm{Hom}(A,A) = \{\mathrm{id}_A\}$. Composition and identity are
trivially defined. This is the ``objects only'' category: all structure has
been stripped away except the objects themselves.

\formalnote{Indiscrete category on $S$}
The \term{indiscrete category} $\mathrm{Indis}(S)$ has $\mathrm{Ob} = S$
and $\mathrm{Hom}(A,B) = \{*_{AB}\}$ for every pair $A, B$ (including
$A = B$, where $*_{AA} = \mathrm{id}_A$). Composition is forced: the only
possible composite of any two composable morphisms is the unique morphism
between the appropriate objects. All axioms hold automatically.

Discrete and indiscrete categories are the two extremes: one has no
morphisms beyond identities, the other has exactly one morphism between
every pair.


% ═══════════════════════════════════════════════════════════════════════════
\section{Isomorphism}
\label{sec:isomorphism}
% ═══════════════════════════════════════════════════════════════════════════

\subsection{Informally: Going There and Back}

Two objects are ``the same for our purposes'' when you can travel from
one to the other and then back, and both round trips (there-and-back, or
back-and-there) leave you exactly where you started --- not approximately
where you started, but \emph{exactly}.

This is the categorical meaning of ``these two things are interchangeable.''
We call this an \term{isomorphism}.

\subsection{Expanding: Round Trips}

\begin{example}{The Two-Road Town}
Suppose two stations, $A$ and $B$, have a road from $A$ to $B$ and a road
from $B$ to $A$. These roads are an isomorphism if: driving from $A$ to $B$
and then immediately back puts you exactly at $A$ with nothing changed, and
driving from $B$ to $A$ and then immediately back puts you exactly at $B$.

This is stronger than just saying the roads exist. It requires that the
round trips are identity journeys.
\end{example}

\begin{example}{A Recipe Step You Can Undo}
Suppose a recipe step ``chill the dough'' has a corresponding step ``bring
dough to room temperature.'' If doing both in either order leaves the dough
in its original state, the two steps are inverse morphisms. Chilling and
then warming = identity. Warming and then chilling = identity.

Note: in reality, chilling and warming might not be perfect inverses
(some moisture is lost, some structure changes). Whether they are isomorphisms
is a question about your specific model of the process, not about the
real-world chemistry.
\end{example}

\begin{dialog}
I see two morphisms $f : A \to B$ and $g : B \to A$. I ask:
\begin{itemize}
  \item Does $g \circ f = \mathrm{id}_A$? (Going from $A$ to $B$ then back
        returns to $A$ as if nothing happened.)
  \item Does $f \circ g = \mathrm{id}_B$? (Going from $B$ to $A$ then back
        returns to $B$ as if nothing happened.)
\end{itemize}
If both hold, $f$ and $g$ are inverse to each other, and both are
isomorphisms.
\end{dialog}

\subsection{Formalizing: The Definition}

A morphism $f : A \to B$ is an \term{isomorphism} if there exists a morphism
$g : B \to A$ such that:
\[
  g \circ f = \mathrm{id}_A \qquad \text{and} \qquad f \circ g = \mathrm{id}_B.
\]
The morphism $g$ is called the \term{inverse} of $f$, written $f^{-1}$.

Two objects $A$ and $B$ are \term{isomorphic}, written $A \cong B$, if there
exists an isomorphism $f : A \to B$.

\textbf{Inverses are unique.} If $g$ and $g'$ are both inverses of $f$, then:
\[
  g = g \circ \mathrm{id}_B = g \circ (f \circ g') = (g \circ f) \circ g'
  = \mathrm{id}_A \circ g' = g'.
\]
Every isomorphism has exactly one inverse.

\subsection{Reorganizing: Which Objects Are Isomorphic?}

\textbf{In $\mathbf{2}$}: the objects $0$ and $1$ are \emph{not} isomorphic.
The only morphism $0 \to 1$ is $\iota$, and there is no morphism $1 \to 0$.
Without a morphism in both directions, there is no isomorphism.

\textbf{In the temperature ordering}: $20°\text{C}$ and $37°\text{C}$ are
not isomorphic. A morphism $20°\text{C} \to 37°\text{C}$ exists (the
``is colder than'' relation), but no morphism goes in the other direction.
In a poset category, two distinct objects are \emph{never} isomorphic.

\textbf{In $\mathbf{1}$}: the single object $*$ is isomorphic to itself,
via $\mathrm{id}_*$. Every object is always isomorphic to itself, via its
identity morphism.

\textbf{In the indiscrete category on $\{A, B\}$}: $A \cong B$. There is a
morphism $*_{AB} : A \to B$ and a morphism $*_{BA} : B \to A$, and their
composites are the unique morphisms $*_{AA} = \mathrm{id}_A$ and
$*_{BB} = \mathrm{id}_B$.

\subsection{Simplifying: Isomorphism in One Line}

$f : A \to B$ is an isomorphism iff $\exists\, g : B \to A$ such that
$g \circ f = \mathrm{id}_A$ and $f \circ g = \mathrm{id}_B$.

The inverse $g$ is then unique and denoted $f^{-1}$.

$A \cong B$ means an isomorphism between them exists.

\subsection{Formally: Properties of Isomorphism}

\formalnote{Isomorphism is an equivalence relation on objects}
\begin{itemize}
  \item \textbf{Reflexive}: $A \cong A$ via $\mathrm{id}_A$ (its own inverse).
  \item \textbf{Symmetric}: if $f : A \xrightarrow{\sim} B$ then
        $f^{-1} : B \xrightarrow{\sim} A$.
  \item \textbf{Transitive}: if $f : A \xrightarrow{\sim} B$ and
        $g : B \xrightarrow{\sim} C$, then $g \circ f : A \xrightarrow{\sim} C$
        with inverse $f^{-1} \circ g^{-1}$.
\end{itemize}

The notation $\xrightarrow{\sim}$ over an arrow marks it as an isomorphism.

\formalnote{Isomorphisms in special categories}
In a thin (preorder) category, the only isomorphisms are identity morphisms
--- since $f : A \to B$ and $g : B \to A$ existing simultaneously forces
$A = B$ (by antisymmetry, if the preorder is a poset). In a groupoid
(a category where \emph{every} morphism is an isomorphism), all morphisms
are invertible.


% ═══════════════════════════════════════════════════════════════════════════
\section{The Opposite Category}
\label{sec:opposite-category}
% ═══════════════════════════════════════════════════════════════════════════

\subsection{Informally: Turning All Arrows Around}

Every category has a mirror image. In the mirror, every arrow points the
other way. An arrow that went from $A$ to $B$ now goes from $B$ to $A$.
Composition still works --- you just have to track the reversal.

This mirror is called the \term{opposite category}.

\subsection{Expanding: Reversing Familiar Examples}

\begin{example}{The Recipe Reversed}
The original recipe category has $f : \mathrm{Chop} \to \mathrm{Sauté}$ and
$g : \mathrm{Sauté} \to \mathrm{Add}$. The opposite category has
$f^{\mathrm{op}} : \mathrm{Sauté} \to \mathrm{Chop}$ and
$g^{\mathrm{op}} : \mathrm{Add} \to \mathrm{Sauté}$.

The original morphism $g \circ f : \mathrm{Chop} \to \mathrm{Add}$ becomes
$(g \circ f)^{\mathrm{op}} : \mathrm{Add} \to \mathrm{Chop}$.

Importantly, $(g \circ f)^{\mathrm{op}} = f^{\mathrm{op}} \circ g^{\mathrm{op}}$.
The order of composition reverses when you reverse arrows. This is the key
bookkeeping fact.
\end{example}

\begin{example}{The Opposite of a Preorder}
The opposite of the temperature preorder has morphisms going from warmer to
colder: $100°\text{C} \to 37°\text{C} \to 20°\text{C} \to 0°\text{C}$.
This is the ``is no colder than'' relation.

In general, the opposite of a preorder category is another preorder category,
with the order reversed.
\end{example}

\begin{dialog}
When I reverse a category, I keep two things straight:
\begin{enumerate}
  \item Every morphism $f : A \to B$ becomes $f^{\mathrm{op}} : B \to A$.
        Domain and codomain swap.
  \item Composition reverses order: if $h = g \circ f$ in $\mathcal{C}$,
        then $h^{\mathrm{op}} = f^{\mathrm{op}} \circ g^{\mathrm{op}}$ in
        $\mathcal{C}^{\mathrm{op}}$.
\end{enumerate}
The objects stay the same. Only the arrows and their order change.
\end{dialog}

\subsection{Formalizing: The Definition}

Given a category $\mathcal{C}$, its \term{opposite category}
$\mathcal{C}^{\mathrm{op}}$ is defined as follows:

\begin{itemize}
  \item $\mathrm{Ob}(\mathcal{C}^{\mathrm{op}}) = \mathrm{Ob}(\mathcal{C})$
        (same objects).
  \item $\mathrm{Hom}_{\mathcal{C}^{\mathrm{op}}}(A, B) =
        \mathrm{Hom}_{\mathcal{C}}(B, A)$ (morphisms are reversed).
  \item Composition in $\mathcal{C}^{\mathrm{op}}$: given
        $f^{\mathrm{op}} \in \mathrm{Hom}_{\mathcal{C}^{\mathrm{op}}}(A, B)$
        and $g^{\mathrm{op}} \in \mathrm{Hom}_{\mathcal{C}^{\mathrm{op}}}(B,C)$,
        their composite is $(f \circ g)^{\mathrm{op}}$, where $\circ$ is
        composition in $\mathcal{C}$.
  \item Identity: $\mathrm{id}_A^{\mathrm{op}} = \mathrm{id}_A$.
\end{itemize}

\textbf{$\mathcal{C}^{\mathrm{op}}$ is a category.} We verify:
\begin{itemize}
  \item \textit{Associativity}: inherited from $\mathcal{C}$, with order
        reversed.
  \item \textit{Unit laws}: $f^{\mathrm{op}} \circ \mathrm{id}_A^{\mathrm{op}}
        = (\mathrm{id}_A \circ f)^{\mathrm{op}} = f^{\mathrm{op}}$, and
        similarly on the other side.
\end{itemize}

\subsection{Reorganizing: The Opposite of the Opposite}

Applying the opposite construction twice returns to the original:
$(\mathcal{C}^{\mathrm{op}})^{\mathrm{op}} = \mathcal{C}$. Reversing arrows
twice gives back the original arrows.

\subsection{Simplifying: Opposite in Notation}

\[
  \mathrm{Hom}_{\mathcal{C}^{\mathrm{op}}}(A, B)
  = \mathrm{Hom}_{\mathcal{C}}(B, A)
\]
\[
  (g \circ_{\mathcal{C}^{\mathrm{op}}} f) = (f \circ_\mathcal{C} g)^{\mathrm{op}}
\]
\[
  (\mathcal{C}^{\mathrm{op}})^{\mathrm{op}} = \mathcal{C}
\]

\subsection{Formally: Why the Opposite Category Matters}

\formalnote{Duality}
The opposite category is the engine of \term{categorical duality}. Any
statement that is true of all categories remains true when you replace every
$\mathcal{C}$ with $\mathcal{C}^{\mathrm{op}}$ --- but the statement it
becomes may sound quite different. Every theorem has a \emph{dual theorem},
obtained by reversing all arrows. This doubles the value of every proof.

\formalnote{Contravariant functors}
In Chapter~3 we will define functors. A \term{contravariant functor} from
$\mathcal{C}$ to $\mathcal{D}$ is simply a (covariant) functor from
$\mathcal{C}^{\mathrm{op}}$ to $\mathcal{D}$. The opposite category allows
us to treat contravariance as a special case of covariance rather than a
separate concept.


% ═══════════════════════════════════════════════════════════════════════════
\section*{Exercises}
\addcontentsline{toc}{section}{Exercises}
% ═══════════════════════════════════════════════════════════════════════════

\begin{exercisesbox}
\begin{qa}
\qa{What four pieces of data specify a category?}{Objects, morphisms, composition, and identities (one per object).}
\qa{What two laws do those pieces obey?}{Associativity of composition, and the unit law for identities.}
\qa{Is ``a set with extra structure'' enough to be a category?}{No. A category is data plus laws; the morphisms and composition are part of the structure.}
\qa{In $\mathbf{0}$, how many morphisms are there?}{Zero. No objects, no morphisms.}
\qa{In $\mathbf{1}$, how many morphisms?}{One: the identity on the unique object.}
\qa{In $\mathbf{2}$, how many morphisms in total?}{Three: two identities and one non-identity arrow.}
\qa{Why does $\mathbf{2}$ have no other arrows?}{There is nothing to compose the existing arrow with that would create a new one.}
\qa{What is a preorder, viewed as a category?}{A category where every hom-set has at most one element.}
\qa{In a preorder-category, when is $A \to B$ inhabited?}{When $A \leq B$ in the preorder.}
\qa{What does antisymmetry give you on top of a preorder?}{Iso-equality: $A \cong B$ implies $A = B$ on the nose.}
\qa{What is a monoid, viewed as a category?}{A category with one object, whose morphisms compose like the monoid operation.}
\qa{In a one-object category, what is the identity morphism?}{The unit of the monoid.}
\qa{What is an isomorphism in a category?}{A morphism with a two-sided inverse.}
\qa{Why must inverses be two-sided?}{Because a one-sided inverse need not be unique or behave well under composition.}
\qa{If $g \circ f = \mathrm{id}_A$ and $f \circ g = \mathrm{id}_B$, what is $g$ called?}{The inverse of $f$ (and vice versa).}
\qa{Are inverses unique when they exist?}{Yes. If $g$ and $g'$ are both inverses of $f$, then $g = g'$.}
\qa{Why is uniqueness of inverses important?}{It lets us speak of \emph{the} inverse, written $f^{-1}$.}
\qa{What does $A \cong B$ mean?}{There exists an isomorphism between $A$ and $B$.}
\qa{Is $\cong$ reflexive? Symmetric? Transitive?}{Yes to all three.}
\qa{What is the opposite category $\mathcal{C}^{\mathrm{op}}$?}{Same objects as $\mathcal{C}$, but every arrow reversed; composition is reversed too.}
\qa{If $f : A \to B$ in $\mathcal{C}$, what is it in $\mathcal{C}^{\mathrm{op}}$?}{An arrow $B \to A$.}
\qa{Why is $(\mathcal{C}^{\mathrm{op}})^{\mathrm{op}} = \mathcal{C}$?}{Reversing arrows twice puts them back the way they were.}
\qa{What is the principle of duality?}{Every theorem about $\mathcal{C}$ yields a theorem about $\mathcal{C}^{\mathrm{op}}$ by reversing arrows.}
\qa{What is a contravariant operation, in opposite-category terms?}{An ordinary (covariant) operation from $\mathcal{C}^{\mathrm{op}}$.}
\qa{Why is a small category called ``small''?}{Its collection of objects (and morphisms) is a set, not a proper class.}
\qa{Is $\mathbf{Set}$ small?}{No. The collection of all sets is too large to be a set.}
\qa{Can a category have a proper class of objects?}{Yes; such categories are called \emph{large}.}
\qa{If $f$ is an iso and $f \circ g$ is defined, can you cancel $f$ from $f \circ g = f \circ h$?}{Yes. Compose with $f^{-1}$ on the left.}
\qa{If $f$ is \emph{not} an iso, can you still cancel?}{Not in general.}
\qa{Where does the cycle ``Informally $\to$ Expanding $\to$ \ldots $\to$ Formally'' end up in a Chapter?}{In the Formal Restatement at the chapter's end.}
\end{qa}
\end{exercisesbox}


% ═══════════════════════════════════════════════════════════════════════════
\section*{Chapter Summary}
\addcontentsline{toc}{section}{Chapter Summary}
% ═══════════════════════════════════════════════════════════════════════════

\begin{summarybox}
A \textbf{category} $\mathcal{C}$ packages: objects, hom-sets, composition,
and identity morphisms, satisfying associativity and the unit laws.

\textbf{Small categories} include $\mathbf{0}$ (empty), $\mathbf{1}$ (one
object), $\mathbf{2}$ (the walking morphism), preorder categories (at most
one morphism between any two objects), and one-object categories (monoids).

An \textbf{isomorphism} $f : A \to B$ has an inverse $g : B \to A$ with
$g \circ f = \mathrm{id}_A$ and $f \circ g = \mathrm{id}_B$. Isomorphic
objects are interchangeable. The inverse is unique.

The \textbf{opposite category} $\mathcal{C}^{\mathrm{op}}$ has the same
objects as $\mathcal{C}$ but all morphisms reversed. Every theorem about
categories has a dual obtained by replacing $\mathcal{C}$ with
$\mathcal{C}^{\mathrm{op}}$.

\medskip
\textbf{The category checklist:}
\begin{enumerate}
  \item Name the objects.
  \item Name the morphisms (with domains and codomains).
  \item Verify composition is defined and associative.
  \item Verify every object has a neutral identity morphism.
\end{enumerate}
\end{summarybox}

% ═══════════════════════════════════════════════════════════════════════════
\section*{Formal Restatement}
\addcontentsline{toc}{section}{Formal Restatement}
% ═══════════════════════════════════════════════════════════════════════════

\begin{formalrestatement}
\textbf{Category.}\quad A category $\mathcal{C}$ consists of
$\mathrm{Ob}(\mathcal{C})$, hom-sets $\mathrm{Hom}(A,B)$ for all $A,B$,
composition $\circ$, and identities $\mathrm{id}_A$, satisfying:
\[
  (h \circ g) \circ f = h \circ (g \circ f),
  \qquad
  f \circ \mathrm{id}_A = f = \mathrm{id}_B \circ f.
\]

\medskip
\textbf{Special classes.}\quad
$\mathcal{C}$ is \emph{small} if $\mathrm{Ob}(\mathcal{C})$ is a set;
\emph{thin} if $|\mathrm{Hom}(A,B)| \leq 1$ for all $A,B$;
\emph{discrete} if $\mathrm{Hom}(A,B) = \varnothing$ for $A \neq B$.

\medskip
\textbf{Isomorphism.}\quad
$f : A \to B$ is an isomorphism iff $\exists\, f^{-1} : B \to A$ with
$f^{-1} \circ f = \mathrm{id}_A$ and $f \circ f^{-1} = \mathrm{id}_B$.
$A \cong B$ means such $f$ exists. The inverse is unique.

\medskip
\textbf{Opposite category.}\quad
$\mathrm{Ob}(\mathcal{C}^{\mathrm{op}}) = \mathrm{Ob}(\mathcal{C})$;
$\mathrm{Hom}_{\mathcal{C}^{\mathrm{op}}}(A,B) = \mathrm{Hom}_\mathcal{C}(B,A)$;
composition $(g^{\mathrm{op}} \circ_{\mathrm{op}} f^{\mathrm{op}}) =
(f \circ g)^{\mathrm{op}}$;
$(\mathcal{C}^{\mathrm{op}})^{\mathrm{op}} = \mathcal{C}$.

\medskip
\noindent\textit{Chapter~3 introduces functors: structure-preserving maps
between categories.}
\end{formalrestatement}
